
\section{Introduction}

Large language models (LLMs) such as GPT-5~\citep{openai2025gpt5}, Gemini-2.5~\citep{comanici2025gemini}, Qwen-3~\citep{qwen3}, and DeepSeek-V3~\citep{deepseekai2024deepseekv3technicalreport} have evolved into the core infrastructure of modern AI. 
Yet as scaling reaches hundreds of billions of parameters, performance gains increasingly depend on a model’s ability to \textbf{reason}—to decompose problems, infer hidden relations, and make consistent multi-step deductions. 
We believe that \textbf{reasoning capability is the essence of intelligence} and the foundation for building \textit{general-purpose agents} that can understand, decide, and act autonomously.

Recent open models highlight this trend. Kimi-K2~\citep{kimiK2}, an open trillion-scale model, focuses primarily on enhancing agentic capability, while DeepSeek-V3~\citep{deepseekai2024deepseekv3technicalreport}, though smaller at 671B parameters, achieves outstanding reasoning performance under efficient sparse scaling. 
{\textbf{Ling~2.0}} is designed to push beyond: \textbf{scaling} a trillion-parameter \textbf{reasoning-oriented} foundation model that maximizes reasoning accuracy and efficiency under sparse activation, establishing a scalable blueprint for next-generation open intelligent systems.

Scaling general reasoning capability to the trillion-parameter level is a central challenge in the evolution of LLMs. 
The key difficulty lies in achieving both \textbf{efficient scaling}—maintaining computational efficiency, stability, and predictability under extreme scale—and \textbf{reasoning enhancement}—ensuring that expanded capacity leads to more consistent and reliable reasoning.

From the scaling perspective, dense architectures incur prohibitive cost, motivating \textbf{high-sparsity designs} that preserve expressiveness while reducing computation. 
Reliable \textbf{scaling prediction} becomes essential to anticipate trillion-scale performance (beyond $1e25$ FLOPs) from smaller-scale experiments. 
In addition, effective \textbf{algorithm-infrastructure co-design} is required to align precision, parallelism, and communication for efficient large-scale execution. 
From the reasoning perspective, maintaining improvement across \textbf{pre-training}, \textbf{mid-training}, and \textbf{post-training} remains difficult. 
Constructing reasoning-centric corpora is resource-intensive, while transferring learned reasoning behaviors across these stages can introduce instability. 
Achieving sustained progress thus requires innovations in both data and training pipeline to balance reasoning accuracy and efficiency. 


To address the intertwined challenges of efficient scaling and sustained reasoning enhancement, Ling 2.0 introduces systematic innovations across four dimensions: model architecture, pre-training, post-training, and infrastructure.


\paratitle{Model Architecture.} 
\begin{itemize}
  \item \textbf{Ling Scaling Laws.} 
  Our unified Ling Scaling Laws, derived from over a \textit{thousand experiments}, guide the hyperparameter and architectural design for trillion-parameter models, ensuring stable and near-optimal training. Crucially, the framework establishes a ``wind tunnel'' for \textit{low-cost, high-fidelity extrapolation} from small-scale trials to trillion-parameter models, cutting validation costs to under 1\% of a full training run and greatly accelerating innovation cycle.
  
  \item \textbf{High-Sparsity MoE with MTP.} 
  Ling 2.0 scales our ``\textit{high-sparsity, fine-grained}'' architecture from 16B to 1T parameters. All models use 256 routed experts, activating 8 experts plus one shared expert per token ($\approx$ 3.5\% activation), realizing \textit{7× efficiency leverage} per the Ling Scaling Law. With aux-loss-free load balancing and MTP, Ling 2.0 maintains high training efficiency while improving logical reasoning, leading to significant math and coding performance gains.
  
\end{itemize}

\paratitle{Pre-Training.} 
\begin{itemize}
  \item \textbf{Reasoning-oriented Data Composition.} 
  Our pre-training corpus prioritizes the \textit{Ling Math} and \textit{Ling Code} datasets, which are tailored for mathematical reasoning and code generation, respectively, yielding a 5-8\% average gain on reasoning benchmarks. Throughout the 20T-token pre-training process, we progressively increase the proportion of reasoning data from 32\% to 46\%, establishing Ling 2.0's inherent reasoning strengths.
    
  \item \textbf{Reasoning Pre-Activation in Mid-Training.} 
  In the mid-training phase, we extend the effective context window and introduce \textit{Chain-of-Thought (CoT) data} to pre-activate reasoning abilities. This strategy raises the ceiling on reasoning performance, and provides a more stable foundation for subsequent fine-tuning and reinforcement learning (RL). 
  
  \item \textbf{Warmup-Stable-Merge (WSM) Scheduler.} 
  To enable a more flexible and effective pre-training process, the Ling 2.0 series adopts the novel WSM (warmup-stable-merge) scheduler, which replaces learning-rate decay with checkpoint merging and delivers 1-2\% average gains across benchmarks. Notably, this advantage persists through subsequent post-training stages. 
\end{itemize}

\paratitle{Post-Training.} 
\begin{itemize}
  \item \textbf{DFT Initialization with Progressive Reasoning Evolution.} Through \textit{Decoupled Fine-Tuning (DFT)} with differentiated system prompts, we establish a diverse, reasoning-focused initialization. Building on this foundation, the \textit{Evolutionary Chain-of-Thought (Evo-CoT)} paradigm progressively deepens reasoning capabilities—enabling Ling 2.0 to surpass state-of-the-art models on competition-level mathematical reasoning benchmark, while requiring 25\% fewer training tokens to reach comparable or better performance.
  \item \textbf{Sentence-Level Policy Optimization.} Introduces \textit{Linguistic-unit Policy Optimization (LPO)}, treating sentences as the fundamental action units for RL updates. This fine-grained optimization strategy shows higher training stability and delivers around 10\% improvements on complex reasoning benchmarks compared to token-level and sequence-level baselines. 
  \item  \textbf{Group-Based Human Preference Alignment.} The \textit{Group Arena Reward (GAR)} mechanism ensures precise intra-group preference alignment in RLHF, better reflecting nuanced human judgments, yielding 2-10\% higher consistency scores in open-ended evaluations.
\end{itemize}

\paratitle{Infrastructure.} 
\begin{itemize}
  \item \textbf{Full-scale FP8 training.} Ling 2.0 represents the largest open-source model trained entirely in FP8 precision. Fine-grained quantization (activations/gradients [1,128]; weights [128,128]) achieves near-lossless accuracy ($\leq$ 0.25 \% gap to BF16 after 900 B tokens) while improving utilization and reducing memory use by over 15 \%.  
  \item \textbf{Heterogeneous fine-grained pipeline.} Interleaved 1F1B scheduling with partial recomputation mitigates pipeline bubbles from heterogeneous modules such as MTP and First-K-Dense, improving throughput by around 40 \%.  
  \item \textbf{Software Engineering for Foundation LLMs.} Guiding a software-engineering-oriented LLMs framework with the 4C (Correct, Consistent, Complete, and Co-Design) principle, incorporating efficient automated iteration, algorithm-system co-design and cross-platform reproducibility to jointly ensures robust trillion-scale development.
\end{itemize}

Based on the above innovations, we release three models of different scales in the Ling 2.0 family:
\begin{itemize}
    \item \texttt{\textbf{Ling-mini-2.0}}: 16B total parameters with 1.4B activated.
    \item \texttt{\textbf{Ling-flash-2.0}}: 103B total parameters with 6.1B activated.
    \item \texttt{\textbf{Ling-1T}}: 1 trillion total parameters with 51B activated.
\end{itemize}


Ling 2.0 is comprehensively evaluated across a wide range of benchmarks spanning mathematics, coding, reasoning, knowledge, alignment, and agentic tasks. The results exhibit a consistent scaling trajectory: as model capacity expands from \texttt{Ling-mini-2.0} to \texttt{Ling-flash-2.0} and \texttt{Ling-1T}, performance across all tasks improve steadily in accordance with the Ling Scaling Law.  

At smaller scales, \texttt{Ling-mini-2.0} achieves performance on par with or exceeding dense models below 10B parameters, while \texttt{Ling-flash-2.0} matches or surpasses dense models below 40B. These findings confirm that Ling 2.0 provides an approximate \textbf{7× efficiency leverage}, delivering dense-level capability with substantially lower active computation.  


At the trillion-parameter scale, \texttt{Ling-1T} establishes a new Pareto frontier of reasoning accuracy versus efficiency, demonstrating “efficient thinking and precise reasoning” on competition-level benchmarks such as AIME 2025. Collectively, these results validate that Ling 2.0 effectively scales reasoning capability with both architectural efficiency and algorithmic alignment, advancing the frontier of open-source language foundation models.

This report focuses on three reflex-grade non-thinking (instruct) models in the Ling 2.0 family—\texttt{Ling-mini-2.0}, \texttt{Ling-flash-2.0}, and \texttt{Ling-1T}. 
These models emphasize general reasoning and instruction-following capability, while the \texttt{\textbf{Ring}} series~\citep{lingteam2025stepevolvesscalingreinforcement}, built upon the same Ling 2.0 base, extends toward deep thinking models.
The remainder of this report introduces the core model architecture, pre-training and post-training methodology, as well as the infrastructure optimizations of Ling 2.0.



